\chapter{Números decimais}\label{ch:g6:decimals}

Entre $2$ e $3$ há um monte de números: $2.5$, $2.71$,
$2.999$\,\dots\ Os \omterm{def:g6:decimals:places}{números decimais} estendem o sistema de \omterm{def:g6:wholes:place}{valor
posicional} para a direita das unidades, e seguem as mesmas regras dos
números naturais --- com algumas armadilhas que este capítulo vai
ensinar você a evitar.

\section{A escrita decimal}

\begin{definition}[Casas decimais]\label{def:g6:decimals:places}
Um \emph{número decimal}\index{número decimal} tem uma parte inteira e
uma parte decimal, separadas pelo \omterm{def:g4:decimals:point}{ponto decimal}. Cada casa à direita do
ponto vale dez vezes menos que a anterior: décimos, centésimos,
milésimos. Por exemplo,
\[
13.407 = 13 + \frac{4}{10} + \frac{0}{100} + \frac{7}{1000}
= 1 \times 10 + 3 + 4 \times \frac{1}{10} + 7 \times \frac{1}{1000}.
\]
\end{definition}

\begin{omfigure}
\begin{tikzpicture}[scale=0.9]
  \draw (0,0) grid[xstep=2.3, ystep=0.85] (11.5,1.7);
  \node[font=\small] at (1.15,1.275) {dezenas};
  \node[font=\small] at (3.45,1.275) {unidades};
  \node[font=\small] at (5.75,1.275) {décimos};
  \node[font=\small] at (8.05,1.275) {centésimos};
  \node[font=\small] at (10.35,1.275) {milésimos};
  \node[font=\large, omDef] at (1.15,0.425) {$1$};
  \node[font=\large, omDef] at (3.45,0.425) {$3$};
  \node[font=\large, omThm] at (5.75,0.425) {$4$};
  \node[font=\large, omThm] at (8.05,0.425) {$0$};
  \node[font=\large, omThm] at (10.35,0.425) {$7$};
  \fill (4.6,0.06) circle (2.5pt);
\end{tikzpicture}

{\small O quadro de \omterm{def:g6:wholes:place}{valor posicional} de $13.407$: o \omterm{def:g4:decimals:point}{ponto decimal} fica
entre as unidades e os décimos.}
\end{omfigure}

\begin{remark}[Zeros que contam e zeros que não contam]
Acrescentar \omterm{def:g1:counting:zero}{zeros} no \emph{fim} da parte decimal não muda nada:
$2.5 = 2.50 = 2.500$. Mas os \omterm{def:g1:counting:zero}{zeros} \emph{entre} algarismos são
essenciais: $13.407 \neq 13.47$. E $0.5$ é bem diferente de $0.05$!
\end{remark}

\begin{method}[Comparar números decimais]\label{met:g6:decimals:compare}
\begin{enumerate}
  \item Compare primeiro as partes inteiras: $7.2 > 5.99$ porque
        $7 > 5$.
  \item Se as partes inteiras forem iguais, compare as partes decimais
        algarismo por algarismo a partir da esquerda --- completar com
        \omterm{def:g1:counting:zero}{zeros} no fim ajuda: para comparar $3.4$ e $3.15$, escreva $3.40$
        e $3.15$; como $40 > 15$ centésimos, $3.4 > 3.15$.
\end{enumerate}
Atenção: \emph{mais comprido não quer dizer maior}. $3.15$ tem mais
algarismos que $3.4$, mas é menor.
\end{method}

\begin{omfigure}
\begin{tikzpicture}[scale=1.1]
  \draw[-Stealth] (-0.3,0) -- (10.6,0);
  \foreach \i in {0,1,...,10} \draw (\i,-0.08) -- (\i,0.08);
  \node[below, font=\small] at (0,-0.1) {$3$};
  \node[below, font=\small] at (5,-0.1) {$3.5$};
  \node[below, font=\small] at (10,-0.1) {$4$};
  \fill[omThm] (1.5,0) circle (2pt) node[above, font=\small] {$3.15$};
  \fill[omDef] (4,0) circle (2pt) node[above, font=\small] {$3.4$};
  \fill[omMeth] (9,0) circle (2pt) node[above, font=\small] {$3.9$};
\end{tikzpicture}

{\small Ampliação da reta numérica entre $3$ e $4$: cada traço é um
décimo. O ponto $3.15$ fica bem no meio, entre $3.1$ e $3.2$.}
\end{omfigure}

\section{Somar, subtrair, multiplicar}

\begin{method}[Contas armadas com decimais]\label{met:g6:decimals:columns}
Para somar ou subtrair \omterm{def:g6:decimals:places}{números decimais}, alinhe os \emph{\omterm{def:g4:decimals:point}{pontos
decimais}} (assim as unidades ficam sob as unidades e os décimos sob os
décimos), completando com \omterm{def:g1:counting:zero}{zeros} no fim se for preciso; depois calcule
como com números naturais e ponha o ponto do resultado sob os outros.
\end{method}

\begin{example}\label{ex:g6:decimals:addsub}
Calcule $13.7 + 2.85$, alinhando os pontos ($13.7 = 13.70$):
\[
\begin{array}{r}
1\,3.7\,0 \\
+\ \ \,2.8\,5 \\
\hline
1\,6.5\,5
\end{array}
\]
Calcule $6 - 2.35$ (escreva $6 = 6.00$): $6.00 - 2.35 = 3.65$.
Confira: $3.65 + 2.35 = 6$.
\end{example}

\begin{example}[Multiplicar decimais]\label{ex:g6:decimals:mult}
Para calcular $2.3 \times 1.4$: multiplique como números naturais,
$23 \times 14 = 322$; depois conte os algarismos decimais dos fatores
($1 + 1 = 2$) e ponha o ponto de modo que o resultado tenha essa
quantidade: $2.3 \times 1.4 = 3.22$. Conferência de tamanho:
$2.3 \times 1.4$ deve dar um pouco mais que $2.3 \times 1 = 2.3$ --- e
$3.22$ dá.
\end{example}

\section{Multiplicar e dividir por 10, 100, 1000}

\begin{proposition}[Deslocar o ponto]\label{prop:g6:decimals:shift}
Multiplicar um \omterm{def:g6:decimals:places}{número decimal} por $10$, $100$, $1000$ desloca o \omterm{def:g4:decimals:point}{ponto
decimal} dele $1$, $2$, $3$ casas para a \emph{direita}; dividir o
desloca para a \emph{esquerda} (completando com \omterm{def:g1:counting:zero}{zeros} quando for
preciso):
\[
3.75 \times 100 = 375, \qquad
42.1 \div 1000 = 0.0421 .
\]
\end{proposition}

\begin{proof}
Multiplicar por $10$ faz cada algarismo valer dez vezes mais: os
décimos viram unidades, as unidades viram dezenas, e assim por diante
--- cada algarismo anda uma coluna para a esquerda no quadro de \omterm{def:g6:wholes:place}{valor
posicional}, o que é o mesmo que deslocar o ponto uma casa para a
direita. Dividir inverte isso.
\end{proof}

\begin{example}[Unidades de medida]\label{ex:g6:decimals:units}
Converter unidades é exatamente esse jogo: como $1$ m $= 100$ cm,
\[
3.75 \text{ m} = 3.75 \times 100 \text{ cm} = 375 \text{ cm},
\qquad
42 \text{ mm} = 42 \div 10 \text{ cm} = 4.2 \text{ cm}.
\]
\end{example}

\section{Arredondamento}

\begin{definition}[Arredondamento]\label{def:g6:decimals:rounding}
O \emph{arredondamento}\index{arredondamento} de um número à unidade
(ou ao décimo, ao centésimo, \dots) é o número mais próximo com essa
precisão. Olhe o algarismo seguinte: se for $0,1,2,3,4$, arredonde para
baixo (mantenha); se for $5,6,7,8,9$, arredonde para cima.
\end{definition}

\begin{example}
$7.38$ arredondado à unidade é $7$ (algarismo seguinte $3$: mantém);
arredondado ao décimo é $7.4$ (algarismo seguinte $8$: sobe). O preço
$4.996$ arredondado ao centésimo é $5.00$ --- o \omterm{def:g6:decimals:rounding}{arredondamento} pode
mudar todos os algarismos!
\end{example}

\section{Exercícios}

\begin{exercise}[$\star$]\label{exo:g6:decimals:1}
Escreva como \omterm{def:g6:decimals:places}{número decimal}: $5 + \dfrac{3}{10} + \dfrac{7}{100}$;
\quad $\dfrac{9}{10}$; \quad $12 + \dfrac{4}{1000}$;
\quad ``oito unidades e cinco centésimos''.
\end{exercise}

\begin{exercise}[$\star$]\label{exo:g6:decimals:2}
Em $86.354$: qual é o algarismo dos décimos? E o dos centésimos? O que
conta o algarismo $8$? Escreva esse número como na
\cref{def:g6:decimals:places}.
\end{exercise}

\begin{exercise}[$\star$]\label{exo:g6:decimals:3}
Copie e complete com $<$, $>$ ou $=$:
\[
5.3 \;?\; 5.29, \qquad
0.7 \;?\; 0.70, \qquad
2.09 \;?\; 2.9, \qquad
14.5 \;?\; 14.49 .
\]
\end{exercise}

\begin{exercise}[$\star$]\label{exo:g6:decimals:4}
Ordene do menor para o maior:
$4.2$; \ $4.05$; \ $4.51$; \ $4.15$; \ $4.5$.
\end{exercise}

\begin{exercise}[$\star$]\label{exo:g6:decimals:5}
Que \omterm{def:g6:decimals:places}{números decimais} correspondem aos pontos $A$, $B$, $C$ numa reta
numérica graduada em décimos, se $A$ está $3$ traços depois do $6$, $B$
está $7$ traços depois do $6$ e $C$ está $2$ traços depois do $7$?
\end{exercise}

\begin{exercise}[$\star$]\label{exo:g6:decimals:6}
Arme as contas e calcule: $45.8 + 7.65$; \quad $23.4 - 8.72$; \quad
$5.6 \times 2.4$ (conte os algarismos decimais!).
\end{exercise}

\begin{exercise}[$\star$]\label{exo:g6:decimals:7}
Calcule sem escrever nada:
\[
6.42 \times 10, \qquad
0.35 \times 1000, \qquad
78.1 \div 100, \qquad
5 \div 1000 .
\]
\end{exercise}

\begin{exercise}[$\star$]\label{exo:g6:decimals:8}
Converta: $2.4$ m em cm; \quad $370$ g em kg; \quad $0.85$ km em m;
\quad $56$ mm em m.
\end{exercise}

\begin{exercise}[$\star$]\label{exo:g6:decimals:9}
Arredonde $23.867$: à unidade; ao décimo; ao centésimo. Arredonde
$9.97$ ao décimo.
\end{exercise}

\begin{exercise}[$\star\star$]\label{exo:g6:decimals:10}
Um pão francês custa $1.15$. Lena compra três pães e paga com uma nota
de $5$. Escreva as duas contas necessárias e dê o troco dela.
\end{exercise}

\begin{exercise}[$\star\star$]\label{exo:g6:decimals:11}
Encontre um \omterm{def:g6:decimals:places}{número decimal} estritamente entre $7.4$ e $7.5$; depois um
entre $3.99$ e $4$. Quantos números desses existem em cada caso?
\end{exercise}

\begin{exercise}[$\star\star\star$]\label{exo:g6:decimals:12}
Usando cada um dos algarismos $2$, $5$, $8$ exatamente uma vez e um
\omterm{def:g4:decimals:point}{ponto decimal}, escreva o maior número possível e depois o menor
possível. (Números como $.58$ não são permitidos: a parte inteira tem
de conter pelo menos um algarismo.)
\end{exercise}

\section{Problema: o número logo depois de 3 não existe}

\begin{problem}[{Problema de fim de semana --- entre dois números
decimais quaisquer sempre há outro: ampliando a reta numérica}]
\label{pb:g6:decimals:1}
Na escada dos números naturais, todo número tem um vizinho de porta:
logo depois de $7$ vem $8$, e nada mora entre os dois. Os \omterm{def:g6:decimals:places}{números
decimais} são um mundo completamente diferente. Qual é o número logo
depois de $3$? É $3.1$? $3.01$? $3.001$? Este problema desenvolve uma
técnica de \emph{ampliação} da reta numérica (continuando o
\cref{exo:g6:decimals:11}) e chega a uma conclusão famosa: o número
logo depois de $3$ \emph{não existe} --- entre dois \omterm{def:g6:decimals:places}{números decimais}
quaisquer, por mais próximos que estejam, sempre há espaço para mais.

\medskip
\noindent\textbf{Parte I --- Ampliando.}
\begin{enumerate}
  \item Qual é maior, $2.999$ ou $3$? Calcule a \omterm{ex:g1:subtraction:difference}{diferença} entre eles.
  \item Trace uma reta numérica de $7.4$ a $7.5$, graduada em
        centésimos, e marque nela $7.42$, $7.45$ e $7.48$. Liste
        \emph{todos} os números com dois algarismos decimais que ficam
        estritamente entre $7.4$ e $7.5$. Quantos são?
  \item Amplie de novo: entre $7.44$ e $7.45$, liste todos os números
        com três algarismos decimais. Quantos são desta vez?
  \item Usando a mesma ideia, explique como contar os números com
        exatamente três algarismos decimais que ficam estritamente
        entre $7.4$ e $7.5$, e dê essa contagem.
  \item Descreva a receita escondida por trás das questões 2 a 4 (a
        \emph{ampliação}): dados dois números que parecem vizinhos,
        como $5.67$ e $5.68$, como é que escrever mais uma casa decimal
        sempre revela um número estritamente entre eles? Aplique a sua
        receita a $5.67$ e $5.68$, e depois a $0.1999$ e $0.2$.
\end{enumerate}

\medskip
\noindent\textbf{Parte II --- O vizinho que não existe.}
\begin{enumerate}[resume]
  \item Tomás afirma: ``o número logo depois de $3$ é $3.1$''. Mostre
        que ele está errado, dando um número estritamente entre $3$ e
        $3.1$. Tomás recua para $3.01$ e depois para $3.001$. Vença
        cada um desses candidatos.
  \item Explique por que \emph{ninguém} pode ganhar esse jogo de você:
        qualquer que seja o número estritamente maior que $3$ proposto
        por Tomás, a receita de ampliação da questão 5 produz um número
        estritamente entre $3$ e a proposta dele. O que isso prova
        sobre ``o número logo depois de $3$''?
  \item Agora o outro lado: Tomás caça o número logo \emph{antes} de
        $3$ e tenta $2.9$, depois $2.99$, depois $2.999$. Vença os três
        candidatos dele. Depois calcule $3 - 2.999$ e descreva (sem
        armar a conta) a \omterm{ex:g1:subtraction:difference}{diferença} entre $3$ e o número escrito com um
        $2$, um ponto e vinte algarismos $9$.
  \item Por que nada disso funciona com os números \emph{naturais}?
        Explique em uma ou duas frases por que não existe número
        natural estritamente entre $7$ e $8$, embora existam muitos
        \omterm{def:g6:decimals:places}{números decimais} ali.
  \item Um comprimento é anunciado como ``$3.7$ cm, arredondado ao
        décimo'' (\cref{def:g6:decimals:rounding}). Dê o menor
        comprimento que arredonda para $3.7$ cm e explique por que um
        comprimento de $3.75$ cm \emph{não} arredonda para $3.7$ ---
        de modo que o comprimento verdadeiro é pelo menos $3.65$ cm e
        estritamente menor que $3.75$ cm.
\end{enumerate}

\medskip
\noindent\textbf{Parte III --- Brincadeiras com algarismos e pontos.}
\begin{enumerate}[resume]
  \item Ordene do menor para o maior: $0.6$; \ $0.58$; \ $0.123$; \
        $0.0999$. Depois explique a um colega, em uma frase, por que
        $0.0999 < 0.6$ mesmo estando escrito com mais algarismos
        (\cref{met:g6:decimals:compare}).
  \item Os preços de uma loja têm sempre exatamente dois algarismos
        decimais. Existe algum \emph{preço} estritamente entre $4.99$ e
        $5$? Compare com a questão 9: o que os preços e os números
        naturais têm em comum?
  \item Usando cada um dos algarismos $2$, $5$, $8$ exatamente uma vez
        e um \omterm{def:g4:decimals:point}{ponto decimal} (parte inteira não vazia, como no
        \cref{exo:g6:decimals:12}), encontre o número mais próximo de
        $6$. Justifique calculando a distância dos seus melhores
        candidatos até $6$.
  \item O número mais famoso da matemática, $\pi$, satisfaz
        $3.141 < \pi < 3.142$. Dê um \omterm{def:g6:decimals:places}{número decimal} estritamente entre
        $3.141$ e $3.142$; depois um estritamente entre $3.1415$ e
        $3.1416$. (Os matemáticos espremem $\pi$ exatamente assim ---
        cada nova casa decimal é mais uma ampliação. Você vai ver $\pi$
        em ação no \cref{ch:g6:measure}.)
  \item O grande final: explique por que há \emph{mais \omterm{def:g6:decimals:places}{números decimais}
        entre $0$ e $1$ do que qualquer número que você consiga
        dizer}. (Quantos números uma ampliação produz? A ampliação pode
        parar algum dia?)
\end{enumerate}
\end{problem}
